\clearpage
\VerdierPhysicalPageAnchor{0176}%
\begingroup\small\setlength{\parskip}{0pt}\setlength{\emergencystretch}{1.5em}
\setlength{\abovedisplayskip}{4pt plus 2pt minus 2pt}%
\setlength{\belowdisplayskip}{4pt plus 2pt minus 2pt}%

\VerdierSemanticAnchor{ch3d-1-2}
\subsection{Les sous-catégories remarquables des catégories dérivées.}

\VerdierSemanticAnchor{ch3d-1-2-1}
\noindent\textbf{Proposition 1.2.1.}\quad
\textit{a) Les sous-catégories
$K^{,*}(\mathcal A)\subset K(\mathcal A)$ sont
$\operatorname{Ac}(\mathcal A)$-localisantes à droite et à gauche
(chap. II, 2.3.5).}

\noindent\textit{b) Les sous-catégories
$K^{+,*}(\mathcal A)\subset K(\mathcal A)$ sont
$\operatorname{Ac}(\mathcal A)$-localisantes à droite.}

\noindent\textit{b)\textsuperscript{bis} Les sous-catégories
$K^{-,*}(\mathcal A)\subset K(\mathcal A)$ sont
$\operatorname{Ac}(\mathcal A)$-localisantes à gauche.}

\noindent\textit{c) La sous-catégorie
$K^{b}(\mathcal A)\subset K^{+}(\mathcal A)$ est
$\operatorname{Ac}^{+}(\mathcal A)$-localisante à gauche.}

\noindent\textit{c)\textsuperscript{bis} La sous-catégorie
$K^{b}(\mathcal A)\subset K^{-}(\mathcal A)$ est
$\operatorname{Ac}^{-}(\mathcal A)$-localisante à droite.}

\noindent\textit{d) Tout objet de $K^{,+}(\mathcal A)$ s'envoie par un
quasi-isomorphisme dans un objet de $K^{+}(\mathcal A)$.}

\noindent\textit{d)\textsuperscript{bis} Tout objet de
$K^{,-}(\mathcal A)$ reçoit par un quasi-isomorphisme un objet de
$K^{-}(\mathcal A)$.}

\noindent\textit{e) Tout objet de $K^{,b}(\mathcal A)$ s'envoie par un
quasi-isomorphisme dans un objet de $K^{+,b}(\mathcal A)$.}

\noindent\textit{e)\textsuperscript{bis} Tout objet de
$K^{,b}(\mathcal A)$ reçoit par un quasi-isomorphisme un objet de
$K^{-,b}(\mathcal A)$.}

\noindent\textit{f) Tout objet de $K^{+,b}(\mathcal A)$ reçoit par un
quasi-isomorphisme un objet de $K^{b}(\mathcal A)$.}

\noindent\textit{f)\textsuperscript{bis} Tout objet de
$K^{-,b}(\mathcal A)$ s'envoie par un quasi-isomorphisme dans un objet de
$K^{b}(\mathcal A)$.}

\VerdierSemanticAnchor{ch3d-1-2-1-truncations}
Pour tout complexe $X$, on désigne par $X(n,+\infty)$ (resp.
$X(-\infty,n)$) le complexe :
\[
X(n,+\infty)=\cdots\longrightarrow0\longrightarrow0
\longrightarrow X^n/\operatorname{Im}(d_X^{n-1})
\longrightarrow X^{n+1}\longrightarrow\cdots
\]
(resp.
\[
X(-\infty,n)=\cdots\longrightarrow X^{n-3}\longrightarrow X^{n-2}
\longrightarrow\operatorname{Ker}(d_X^{n-1})
\longrightarrow0\longrightarrow0\longrightarrow\cdots
\])
et par $p\colon X\longrightarrow X(n,+\infty)$ (resp.
$i\colon X(-\infty,n)\longrightarrow X$) le morphisme de complexes :
\[
\begin{aligned}
p^q&=\operatorname{identité},&&q>n,\\
p^n&\colon X^n\longrightarrow X^n/\operatorname{Im}(d_X^{n-1})
  &&\text{(morphisme canonique)},\\
p^q&=0,&&q\leq n-1
\end{aligned}
\]
\endgroup
